\documentclass[11pt]{article}
\usepackage[margin=1in]{geometry}
\usepackage{amsmath,amssymb,amsthm}
\usepackage{booktabs}
\usepackage{hyperref}

\newtheorem{theorem}{Theorem}
\newtheorem{lemma}[theorem]{Lemma}
\newtheorem{proposition}[theorem]{Proposition}
\newtheorem{corollary}[theorem]{Corollary}
\theoremstyle{remark}
\newtheorem{remark}[theorem]{Remark}

\newcommand{\NSym}{\mathbf{Sym}}
\newcommand{\Prim}{\operatorname{Prim}}

\title{Disproof of Conjecture 00000000438:\\
the degree-3 primitives of $\NSym$ form a $2$-dimensional space, but $A(3,1)\in\{1,4\}$}
\author{jilint777}
\date{2026-10-04}

\begin{document}
\sloppy
\maketitle

\begin{abstract}
Conjecture 00000000438 asserts that the primitive Lie algebra of the Hopf
algebra $\NSym$ of noncommutative symmetric functions has $n$-th homogeneous
dimension equal to the Eulerian number $A(n,1)$. We show that the degree-$3$
primitive space $\Prim(\NSym)_3$ has dimension exactly $2$, with basis
$\Psi_3$ and $[S_2,S_1]$. The Eulerian number $A(3,1)$, however, is $4$ in the
convention ``permutations with one descent'' and $1$ in the $1$-based
convention of OEIS A008292. In fact $2$ is not of the form $A(m,1)$ for any
$m$ in either convention, so no shift of the index can save the
formula. A self-contained Lean~4 development (core library only) builds
$\NSym$ in degree $3$ from its defining coproduct, proves that the primitive space
has dimension $2$, defines the Eulerian numbers by counting descents and
ascents of permutations, and proves that the dimension clause is false. More
generally the dimensions are $1,1,2,3,6,9,18,30,56,99$ for $n=1,\dots,10$, which
match $A(n,1)$ only at $n=2$ (both conventions) and at $n=1$ (1-based convention).
\end{abstract}

\section{The conjecture and the clause we refute}

The statement reads:
\begin{quote}
\emph{Definition:} The copying structure of noncommutative symmetric
functions Sym. \emph{Conjecture:} As a coalgebra, the primitive Lie algebra of
Sym is isomorphic to the generalized Lie $n$-algebra, whose $n$-th homogeneous
dimension equals the Eulerian number $A(n,1)$ (verifiable for $n\le10$).
\end{quote}
The Chinese version says the same. The ``copying structure'' is the coproduct
(comultiplication) of $\NSym$, and its ``primitive Lie algebra'' is
$\Prim(\NSym)=\{x:\Delta x=x\otimes1+1\otimes x\}$, a graded Lie algebra
under the commutator. The phrase ``generalized Lie $n$-algebra'' names no
standard object and is not defined in the statement. The only checkable
content is the dimension formula, and it is meant to be checked on the
primitive Lie algebra itself (``verifiable for $n\le10$''). An isomorphism of graded
objects preserves homogeneous dimensions, so the conjecture implies
\begin{equation}\label{eq:clause}
\dim \Prim(\NSym)_n \;=\; A(n,1)\qquad\text{for all } n\ge1 .
\end{equation}
We refute \eqref{eq:clause}. If the conjecture is read as a conjunction
(``$\Prim(\NSym)\cong L$'' and ``$\dim L_n=A(n,1)$''), a graded isomorphism
transfers the second conjunct to $\Prim(\NSym)$, so refuting
\eqref{eq:clause} refutes the conjecture. The isomorphism part is therefore never
needed.

\paragraph{Readings covered.}
\begin{itemize}
\item \emph{Eulerian convention.} $A(n,k)$ is the number of permutations of
$\{1,\dots,n\}$ with $k$ descents (Graham--Knuth--Patashnik), giving
$A(3,1)=4$. In the $1$-based convention (OEIS A008292, $T(n,k)$ = permutations with $k$
ascending runs, i.e.\ $k-1$ descents), $A(3,1)=1$. Ascents give the same
numbers as descents. Section~\ref{sec:euler} shows that $A(m,1)\in\{2^m-m-1,\;1\}$
is never $2$. So $\dim\Prim(\NSym)_3=2$ refutes \eqref{eq:clause} under every
convention and every index shift $n\mapsto n+s$.
\item \emph{Coefficients.} We work over $\mathbb Q$ (any field of characteristic $0$),
the standard setting. The degree-$3$ computation is valid over every commutative
ring: there $\Prim(\NSym)_3$ is free of rank $2$ (Remark~\ref{rem:ring}).
\item \emph{Which algebra.} The statement explicitly says \emph{noncommutative}
symmetric functions, i.e.\ $\NSym$ of Gelfand, Krob, Lascoux, Leclerc, Retakh and Thibon~\cite{GKLLRT}.
For the commutative ring of symmetric functions, or for QSym, every primitive space
is $1$-dimensional (spanned by $p_n$, respectively $M_{(n)}$). That would agree
with the degenerate $1$-based value $A(n,1)=1$, but it is not the algebra named
in the conjecture.
\end{itemize}

\section{\texorpdfstring{The Hopf algebra $\NSym$}{The Hopf algebra NSym}}

$\NSym=\mathbb Q\langle S_1,S_2,\dots\rangle$ is the free associative algebra on
noncommuting generators $S_k$ of degree $k$. Its degree-$n$ part $\NSym_n$ has
basis $S^I=S_{i_1}S_{i_2}\cdots S_{i_\ell}$, where $I=(i_1,\dots,i_\ell)$ runs over the
$2^{n-1}$ compositions of $n$. The coproduct is the algebra morphism
$\Delta:\NSym\to\NSym\otimes\NSym$ with
\[
\Delta S_k=\sum_{i+j=k}S_i\otimes S_j\qquad(S_0=1),
\]
so $\Delta S^I=\Delta S_{i_1}\cdots\Delta S_{i_\ell}$ \cite[\S3]{GKLLRT}. The power sums of the
first kind $\Psi_k$ are defined by Newton's formula
$kS_k=\sum_{j=0}^{k-1}S_j\Psi_{k-j}$. They are primitive, and $\Prim(\NSym)$ is the free
Lie algebra generated by $\Psi_1,\Psi_2,\dots$ \cite{GKLLRT}. We do not need
this structure theorem for degree $3$, which we compute from scratch.

Write $\bar\Delta x=\Delta x-x\otimes1-1\otimes x$ for $x$ of positive degree. Then
$x$ is primitive iff $\bar\Delta x=0$, and $\bar\Delta$ is linear.

\section{\texorpdfstring{The degree-3 primitives}{The degree-3 primitives}}

\begin{lemma}\label{lem:coprod}
In degree $3$ (basis $S_3,S_{21},S_{12},S_{111}$, where $S_{21}=S_2S_1$ etc.):
\begin{align*}
\bar\Delta S_3&=S_2\otimes S_1+S_1\otimes S_2,\\
\bar\Delta S_{21}=\bar\Delta S_{12}&=S_2\otimes S_1+S_1\otimes S_2+S_{11}\otimes S_1+S_1\otimes S_{11},\\
\bar\Delta S_{111}&=3\,S_{11}\otimes S_1+3\,S_1\otimes S_{11}.
\end{align*}
\end{lemma}
\begin{proof}
Expand the products. For example,
$\Delta S_{21}=(S_2\otimes1+S_1\otimes S_1+1\otimes S_2)(S_1\otimes1+1\otimes S_1)
=S_{21}\otimes1+S_2\otimes S_1+S_{11}\otimes S_1+S_1\otimes S_{11}+S_1\otimes S_2+1\otimes S_{21}$.
$\Delta S_{12}$ gives the same six terms, with $S_{12}$ in place of $S_{21}$ in the two outer ones.
$\Delta S_{111}=(S_1\otimes1+1\otimes S_1)^3$ has middle terms $3S_{11}\otimes S_1+3S_1\otimes S_{11}$.
\end{proof}

\begin{theorem}\label{thm:main}
$x=aS_3+bS_{21}+cS_{12}+dS_{111}$ is primitive iff
\[
a+b+c=0\qquad\text{and}\qquad b+c+3d=0 .
\]
Hence $\Prim(\NSym)_3=\{d\,(3,0,-3,1)+b\,(0,1,-1,0)\}$, and
$\dim_{\mathbb Q}\Prim(\NSym)_3=2$. A basis is
\[
\Psi_3=3S_3-S_{21}-2S_{12}+S_{111},\qquad [S_2,S_1]=S_{21}-S_{12}.
\]
\end{theorem}
\begin{proof}
By Lemma~\ref{lem:coprod},
$\bar\Delta x=(a+b+c)(S_2\otimes S_1+S_1\otimes S_2)+(b+c+3d)(S_{11}\otimes S_1+S_1\otimes S_{11})$.
The four tensors $S_J\otimes S_K$ are distinct basis elements of $\NSym\otimes\NSym$,
so $\bar\Delta x=0$ iff both brackets vanish. The two equations have
rank $2$, with solutions $a=3d$, $c=-b-3d$ for free $b,d$, so the kernel has dimension
$4-2=2$. $\Psi_3$ ($a,b,c,d=3,-1,-2,1$) and $[S_2,S_1]$ ($0,1,-1,0$) both satisfy the
equations. They are independent: the $S_3$-coordinate forces the coefficient of
$\Psi_3$ to be $0$, and then the $S_{21}$-coordinate forces the other one to be $0$.
The expression for $\Psi_3$ follows from Newton's formula: $\Psi_1=S_1$, $\Psi_2=2S_2-S_{11}$,
$3S_3=\Psi_3+S_1\Psi_2+S_2\Psi_1=\Psi_3+2S_{12}-S_{111}+S_{21}$.
\end{proof}

\begin{remark}\label{rem:ring}
The proof uses no division. Over any commutative ring $R$ the module $\Prim(\NSym_R)_3$
is free with basis $\{(3,0,-3,1),(0,1,-1,0)\}$. In particular its dimension is $2$
over every field.
\end{remark}

\section{Eulerian numbers}\label{sec:euler}

\begin{lemma}\label{lem:euler}
For $m\ge1$, the number of permutations of $\{1,\dots,m\}$ with exactly one descent is
$2^m-m-1$. Exactly one permutation (the identity) has no descent.
By the reversal $w\mapsto w_m\cdots w_1$, the same holds for ascents.
\end{lemma}
\begin{proof}
A permutation with at most one descent is an increasing word on a set $T$ followed by
an increasing word on its complement. It is determined by $T$, giving $2^m$ choices.
It has no descent iff $T=\{1,\dots,k\}$ for some $0\le k\le m$, which gives $m+1$
choices that all produce the identity. Every other $T$ gives a distinct permutation
with exactly one descent.
\end{proof}

So $A(m,1)=2^m-m-1$ in the descent convention, giving $0,1,4,11,26,57,120,247,502,1013$
for $m=1,\dots,10$, and $A(m,1)=1$ for all $m\ge1$ in the $1$-based convention.
Since $2^{m+1}-(m+1)-1-(2^m-m-1)=2^m-1\ge0$, the sequence $2^m-m-1$ is
nondecreasing. It takes the values $0,0,1,4$ at $m=0,1,2,3$, so it never equals $2$. Thus

\begin{corollary}\label{cor:main}
$\dim\Prim(\NSym)_3=2\neq A(m,1)$ for every $m\ge0$ and every convention above.
In particular \eqref{eq:clause} fails at $n=3$: $2\ne4$ and $2\ne1$.
\end{corollary}

\section{All degrees up to 10}

$\NSym$ is the universal enveloping algebra of the free Lie algebra $\Prim(\NSym)$
generated by $\Psi_1,\Psi_2,\dots$. By Poincar\'e--Birkhoff--Witt,
$\prod_{k\ge1}(1-t^k)^{-p_k}=\sum_n\dim\NSym_n\,t^n=\frac{1-t}{1-2t}$, where
$p_k=\dim\Prim(\NSym)_k$. Solving for $p_k$ gives the first row of the table below. The
same values come from a direct kernel computation and from duality with QSym
(Section~\ref{sec:verify}).

\begin{center}
\begin{tabular}{lcccccccccc}
\toprule
$n$ & 1&2&3&4&5&6&7&8&9&10\\
\midrule
$\dim\Prim(\NSym)_n$ & 1&1&2&3&6&9&18&30&56&99\\
$A(n,1)$, one descent & 0&1&4&11&26&57&120&247&502&1013\\
$A(n,1)$, 1-based & 1&1&1&1&1&1&1&1&1&1\\
\bottomrule
\end{tabular}
\end{center}
Agreement with both conventions holds only at $n=2$. At $n=1$ only the $1$-based value agrees.
From $n=3$ on, the formula fails in every degree up to $10$ under both conventions.
Only $n=3$ is needed for the disproof.

\section{Lean formalization}

\texttt{lean4/Main.lean} (Lean 4.19.0, core library only, no Mathlib) contains:
\begin{itemize}
\item \emph{Objects.} Compositions are lists of naturals, and \texttt{comps n} lists those of $n$.
\texttt{comps\_valid} checks that it lists $2^{n-1}$ distinct compositions of $n$ for $n\le8$.
A basis tensor $S^J\otimes S^K$ is the pair $(J,K)$, and a tensor is a formal integer
combination of such pairs, compared coefficientwise (\texttt{coeffT}).
\texttt{deltaS I} computes $\Delta S^I$ as the product of the $\Delta S_{i}$ in the
tensor-square algebra. An element $\sum_I x(I)S^I$ of degree $n$ is given by its coefficient
function. \texttt{IsPrimitive n x} says that every coefficient of
$\sum_I x(I)(\Delta S^I-S^I\otimes1-1\otimes S^I)$ vanishes. \texttt{PrimDim n m} says
that there are $m$ linearly independent primitives and that every $m+1$ primitives
are linearly dependent, which is the definition of dimension. Coefficients are integers.
Linear (in)dependence over $\mathbb Z$ and over $\mathbb Q$ coincide (clear denominators),
and rational primitives are rational multiples of integral ones.
\item \texttt{isPrimitive\_three\_iff}: Theorem~\ref{thm:main}'s criterion, for all $x$.
It uses a support lemma: only finitely many keys can carry nonzero coefficients.
The coefficient table on those keys is evaluated by \texttt{decide}.
\item \texttt{primDim\_three : PrimDim 3 2}. The witnesses are $\Psi_3$ and $[S_2,S_1]$.
For the upper bound, any three primitives satisfy an explicit integer relation built from
$2\times2$ minors (\texttt{dep3}, \texttt{three\_dep}).
\item \texttt{primDim\_three\_unique}: \texttt{PrimDim 3 m} implies $m=2$.
\item Eulerian numbers: \texttt{perms n} enumerates permutations of $[1,\dots,n]$. For $n\le5$
it lists $n!$ distinct permutations, and for $n=3$ the membership criterion is checked.
\texttt{eulerDes0}, \texttt{eulerDes1}, \texttt{eulerAsc0} and \texttt{eulerAsc1}
count permutations by descents or ascents, $0$- or $1$-based.
\texttt{eulerian\_three} gives $A(3,1)=4,1,4,1$.
\texttt{eulerian\_ne\_two} shows that none of the four equals $2$ at $m\le6$.
\item \texttt{DimClause A}: $\forall n\ge1$, \texttt{PrimDim n (A n 1)}.
\texttt{conjecture\_00000000438\_false} proves $\neg$\texttt{DimClause} for all four
conventions. \texttt{conjecture\_00000000438\_false\_degree3} packages
$\dim\Prim(\NSym)_3=2$ (existence and uniqueness) and $\dim\ne A(m,1)$ for every
$m\le6$ and every convention.
\end{itemize}
There is no \texttt{sorry}, \texttt{native\_decide} or added axiom. \texttt{\#print axioms}
reports at most \texttt{propext}, \texttt{Classical.choice} and \texttt{Quot.sound}.

\paragraph{Not in Lean.} The statement $A(m,1)\ne2$ for $m\ge7$ (Lemma~\ref{lem:euler}),
the dimensions in degrees $n\ne3$, and the equality of $\mathbb Z$- and $\mathbb Q$-dimensions are
proved here and checked by \texttt{verify.py}. None of them is needed for the
refutation at $n=3$ under the stated conventions. The undefined
``generalized Lie $n$-algebra'' is not formalized and is not needed.

\section{Independent verification}\label{sec:verify}

\texttt{verify.py} uses only the Python 3 standard library. It computes
$\dim\Prim(\NSym)_n$ in three ways:
\begin{enumerate}
\item as the kernel of $\bar\Delta$, by exact rational elimination for $n\le8$ and modulo
the prime $2^{61}-1$ for $n\le10$;
\item by duality with QSym, as $2^{n-1}$ minus the rank of the degree-$n$ quasi-shuffle
products $M_\alpha M_\beta$, for $n\le7$;
\item from the PBW identity above, for $n\le10$.
\end{enumerate}
All three agree. It computes Eulerian numbers by brute force over all permutations
(descents and ascents, $n\le9$) and by the closed formula ($n\le10$), and it prints
the degree-$3$ computation in full.

\begin{thebibliography}{9}
\bibitem{GKLLRT} I.~M.~Gelfand, D.~Krob, A.~Lascoux, B.~Leclerc, V.~S.~Retakh, J.-Y.~Thibon,
\emph{Noncommutative symmetric functions}, Adv.\ Math.\ 112 (1995), 218--348.
\bibitem{GKP} R.~L.~Graham, D.~E.~Knuth, O.~Patashnik, \emph{Concrete Mathematics}, 2nd ed.,
Addison--Wesley, 1994, \S6.2.
\bibitem{OEIS} OEIS Foundation, \emph{The On-Line Encyclopedia of Integer Sequences},
A008292 (Eulerian numbers) and A000295 ($2^n-n-1$).
\end{thebibliography}

\end{document}
